% GitHub Repo and Documentation: https://github.com/celiobjunior/resume-template
% Copyright © 2025 Celio B Junior. All rights reserved.
% 
% Licensed under the Apache License, Version 2.0 (the "License");
% you may not use this file except in compliance with the License.
% You may obtain a copy of the License at
%
%     http://www.apache.org/licenses/LICENSE-2.0
%
% This template follows best practices from README.md

% Início do documento LaTeX: define tipo e formato do documento
% a4paper = tamanho A4, 10pt = tamanho base da fonte
\documentclass[a4paper,10pt]{article}

% --- PACOTES: BASE / IDIOMA ---
\usepackage[utf8]{inputenc}
\usepackage[T1]{fontenc}
\usepackage[portuguese]{babel}

% --- PACOTES: LAYOUT / TIPOGRAFIA ---
\usepackage{geometry}
\usepackage{parskip}
\usepackage{microtype}
\usepackage{enumitem}
\usepackage{titlesec}
\usepackage{array}
\usepackage{tabularx}

% --- PACOTES: LINKS / BOOKMARKS ---
\usepackage{hyperref}
\usepackage{bookmark}
\usepackage{xurl}

% --- CONFIGURAÇÃO: MARGENS / ESPAÇAMENTO ---
\geometry{top=1.5cm, bottom=1.5cm, left=1.5cm, right=1.5cm} % Margens da página
\setcounter{secnumdepth}{0}                                 % Remove numeração de seção
\setlist[itemize]{
    leftmargin=0.75em, % Recuo da lista (mais perto da margem da página)
    itemsep=0.2em,     % Espaço entre itens
    topsep=0.25em,     % Espaço antes/depois da lista
    parsep=0em,        % Espaço entre parágrafos no mesmo item
    partopsep=0em      % Espaço extra ao iniciar a lista
}

% Remove números de página e cabeçalhos para visual limpo do currículo
\pagestyle{empty}

% --- CONFIGURAÇÃO: METADADOS / LINKS ---
\hypersetup{
    pdftitle={Karolayne Amábile Brito Borges},   % Título do PDF
    pdfauthor={Karolayne Amábile Brito Borges},     % Autor do PDF
    colorlinks=true,          % Usa cor nos links (sem caixa)
    linkcolor=black,          % Cor de links internos
    urlcolor=black,           % Cor de URLs
    citecolor=black,          % Cor de citações
    bookmarksdepth=2          % Barra lateral: seção e subseção
}

% --- CONFIGURAÇÃO: TÍTULOS ---
\titleformat{\section}
{\Large\bfseries}
{}
{0em}
{}
[\titlerule\vspace{0.5ex}]
\titleformat{\subsection}
{\normalsize\bfseries}
{}
{0em}
{}
\titlespacing*{\subsection}{0pt}{0.35em}{0.15em}

% Macro para entradas do currículo com bookmark no PDF
\newcounter{cventry}
\newcommand{\cventry}[4]{%
    \refstepcounter{cventry}%
    \phantomsection%
    \pdfbookmark[2]{#1}{cventry-\thecventry}%
    \noindent\begin{tabularx}{\textwidth}{@{}>{\raggedright\arraybackslash}X >{\raggedleft\arraybackslash}X@{}}
    \textbf{#1} & #2 \\
    \textit{#3} & \textit{#4} \\
    \end{tabularx}
}

% --- INÍCIO DO DOCUMENTO ---
\begin{document}

% --- CABEÇALHO ---
% Substitua por suas informações pessoais
\begin{center}
    {\LARGE \textbf{Karolayne Amábile Brito Borges}} 
    \\ [0.1cm]
    Anápolis, Goiás
    {\textbullet}
    \href{mailto:karolayneamabile@gmail.com}{karolayneamabile@gmail.com} 
    {\textbullet}
    \href{https://linkedin.com/in/karolayneamabile}{linkedin.com/in/karolayneamabile}
    {\textbullet}
    \href{https://github.com/KarolayneAmabile}{github.com/KarolayneAmabile}
\end{center}

% --- SEÇÕES ---

\section{Resumo}

DevOps Engineer Jr com experiência em infraestrutura cloud-native, Kubernetes, automação e práticas GitOps. Atuo com IaC, pipelines CI/CD, observabilidade e gerenciamento de secrets em ambientes Kubernetes. AWS Certified Cloud Practitioner e graduanda em Ciência da Computação pelo IFG.

% ------

\section{Experiência}
    \cventry{Soliton}{Remoto}{DevOps Engineer Jr}{Abril, 2025 - Maio, 2026}
        \begin{itemize}
            \item Reduzi o tempo de deploy da plataforma OpenEDX em 5x ao substituir a 
            camada de abstração do Tutor por uma pipeline customizada com GitHub Actions 
            (cache otimizado + Dockerfile direto), ArgoCD e ArgoCD Image Updater.
            \item Projetei e administrei clusters Kubernetes multi-purpose (segurança, 
            monitoramento, operações e produção), garantindo isolamento e governança 
            por ambiente.
            \item Desenvolvi módulos Terraform/OpenTOFU e Helm Charts para 
            padronização e automação de infraestrutura em ambientes Kubernetes, reduzindo 
            inconsistências entre deploys e garantindo o padrão GitOps.
            \item Implementei uma stack de observabilidade com Grafana, Prometheus, Loki e 
            Alloy, habilitando monitoramento e alertas para ambientes cloud-native.
            \item Automatizei o gerenciamento de secrets com Vault e External Secrets 
            Operator, eliminando credenciais hardcoded nos repositórios.
            \item Desenvolvi API REST em Rust seguindo arquitetura hexagonal, garantindo
            baixo acoplamento entre domínio, aplicação e infraestrutura.
        \end{itemize}

% ------

\section{Certificações}
    \begin{itemize}
        \item \textbf{AWS Certified Cloud Practitioner} — Amazon Web Services
    \end{itemize}

% ------

\section{Habilidades}
    \begin{itemize}
        \item \textbf{Cloud \& IaC:} AWS, Terraform/OpenTOFU, Terragrunt, Helmfile
        \item \textbf{Kubernetes \& GitOps:} Kubernetes (EKS, k0s), ArgoCD, Helm, 
        Cilium, Cloudflared, cert-manager
        \item \textbf{CI/CD:} GitHub Actions, ArgoCD, ArgoCD Image Updater
        \item \textbf{Observabilidade:} Grafana, Prometheus, Loki, Alloy
        \item \textbf{Storage \& Backup:} Longhorn, Rook-Ceph, Velero
        \item \textbf{Segurança \& Secrets:} Vault, Vaultwarden, External Secrets Operator
        \item \textbf{Linguagens:} Rust, YAML
        \item \textbf{Idiomas:} Português (Nativo), Inglês (Intermediário)
    \end{itemize}

% ------

\section{Educação}
    \cventry{Instituto Federal de Goiás (IFG)}{Goiás, Brasil}
    {Bacharelado em Ciência da Computação}{2024 - 2027}
    \begin{itemize}
        \item IRA: 8,0 — cursando 5º período
        \item Disciplinas relevantes: Arquitetura de Computadores, Sistemas Operacionais,
        Redes de Computadores, Engenharia de Software, Estrutura de Dados, Bancos de Dados
    \end{itemize}
\end{document}